\noindent\textbf{Elección del procedimiento.} La componente vertical de la rotación terrestre es $\Omega\sin\varphi$. El término horizontal de Coriolis utiliza su doble: $f=2\Omega\sin\varphi$. Convertir latitudes de grados a radianes antes de evaluar el seno.

\begin{enumerate}

\item Sustituir la latitud de Kurú.

\[\fitmath{f=2\cdot \Omega\cdot \sin\left(\varphi\right)}\]

\[\fitmath{f=2\cdot (7.292\times10^{-5})\cdot \sin\left((0.090757)\right)\approx 1.321784\times10^{-5}\,\mathrm{1/s}}\]

\item Usar la misma rotación para la comparación con León.

\[\fitmath{\mathrm{fLeon}=2\cdot \Omega\cdot \sin\left(\mathrm{phiLeon}\right)}\]

\[\fitmath{\mathrm{fLeon}=2\cdot (7.292\times10^{-5})\cdot \sin\left((0.7435102613)\right)\approx 9.871559\times10^{-5}\,\mathrm{1/s}}\]

\item Comparar mediante un cociente adimensional.

\[\fitmath{\mathrm{ratio}=\frac{f}{\mathrm{fLeon}}}\]

\[\fitmath{\mathrm{ratio}=\frac{(1.321783791\times10^{-5})}{(9.871559142\times10^{-5})}\approx 0.1338982\,}\]

\end{enumerate}
